\documentclass[conference]{IEEEtran}
\IEEEoverridecommandlockouts
\usepackage[utf8]{inputenc}
\usepackage[T5]{fontenc}
\usepackage{amsmath,amssymb}
\usepackage{graphicx}
\usepackage{booktabs}
\usepackage{multirow}
\usepackage{hyperref}
\usepackage{xcolor}

\title{Tone-Aware Multi-Task Fine-Tuning for Noise-Robust Vietnamese Automatic Speech Recognition}

\author{%
\IEEEauthorblockN{Author 1, Author 2, Author 3, Author 4}
\IEEEauthorblockA{FPT University, Ho Chi Minh City, Vietnam\\
Email: \{...
\}@fpt.edu.vn}
}

\begin{document}
\maketitle

\begin{abstract}
Vietnamese automatic speech recognition remains difficult under real-world noise because lexical tone and diacritics are vulnerable to acoustic degradation. This paper investigates whether decoder-side tone supervision improves noise robustness when adapting Whisper/PhoWhisper with parameter-efficient LoRA fine-tuning. We construct a reproducible noisy Vietnamese ASR benchmark by mixing clean speech with public noise sources at controlled SNR levels, and evaluate standard WER/CER together with Vietnamese-specific tone, diacritic, final-consonant, and short-word deletion metrics. Experiments compare zero-shot Whisper/PhoWhisper, clean LoRA, noisy LoRA, tone-aware multi-task LoRA, and enhancement-based preprocessing. Results show ...
\end{abstract}

\begin{IEEEkeywords}
Vietnamese ASR, Whisper, PhoWhisper, LoRA, tonal languages, noise robustness, multi-task learning.
\end{IEEEkeywords}

\section{Introduction}
Vietnamese is a tonal language, so ASR mistakes involving tone marks and diacritics may change lexical meaning even when the base character sequence is partially correct. Recent pretrained ASR models such as Whisper and PhoWhisper provide strong baselines, but their robustness under realistic Vietnamese noise and their tone-specific error behavior remain insufficiently understood.

\textbf{Research question.} Does explicit tone supervision during parameter-efficient fine-tuning improve Vietnamese ASR robustness under controlled real-world noise?

\textbf{Contributions.}
\begin{itemize}
    \item A reproducible noisy Vietnamese ASR benchmark protocol with fixed noise types, SNR levels, and manifest release.
    \item A compute-aware LoRA adaptation framework for PhoWhisper with decoder-side tone supervision.
    \item Vietnamese-specific evaluation metrics for tone, diacritic, final consonant, and short-word errors.
    \item An analysis of whether speech enhancement improves WER/CER while harming or preserving tone-related metrics.
\end{itemize}

\section{Related Work}
\subsection{Vietnamese ASR}
Summarize VIVOS, Common Voice, FLEURS, PhoWhisper, VietSuperSpeech, PhoASR/Vietnamese ASR Revisit, and VLSP-style evaluations. Emphasize dataset type: read speech, crowd-sourced speech, conversational speech, pseudo-labeled speech.

\subsection{Noise-Robust ASR}
Discuss waveform noise mixing, SpecAugment, robust self-supervised representations, and speech enhancement. Clearly distinguish training-time augmentation from test-time enhancement.

\subsection{Tone and Diacritic Modeling}
Discuss tonal-language ASR, pitch/F0 features, tone recognition, Vietnamese diacritic restoration, and why post-processing alone is not equivalent to acoustic tone preservation.

\subsection{Parameter-Efficient Fine-Tuning}
Discuss LoRA and why PEFT is necessary under 8GB GPU constraints.

\section{Method}
\subsection{Noisy Benchmark Construction}
Let $x$ denote clean speech and $n$ denote a sampled noise segment. For target SNR $s$, the noisy signal is
\begin{equation}
    \tilde{x} = x + \alpha n, \quad
    \alpha = \sqrt{\frac{P_x}{10^{s/10} P_n}},
\end{equation}
where $P_x$ and $P_n$ are average signal powers.

\subsection{Tone Label Extraction}
Vietnamese syllables are mapped to six tone classes: ngang, sắc, huyền, hỏi, ngã, nặng. Tokens with digits, foreign words, acronyms, or ambiguous normalization are masked with $m_i=0$.

\subsection{Tone-Aware Multi-Task Objective}
The ASR model is optimized with
\begin{equation}
    \mathcal{L}=\mathcal{L}_{ASR}+\lambda \mathcal{L}_{tone},
\end{equation}
where
\begin{equation}
    \mathcal{L}_{tone}= -\frac{1}{N_{valid}} \sum_i m_i \sum_c y_{i,c}\log \hat{y}_{i,c}.
\end{equation}

\subsection{Token--Syllable Alignment}
Because Whisper uses subword tokens, we assign each syllable-level tone target to either all subtokens or the final subtoken of the syllable, and mask special tokens. Both policies are treated as ablations.

\section{Experiments}
\subsection{Datasets}
Report train/dev/test hours, speakers, utterances, average duration, and license. Separate clean train data from noisy evaluation data.

\subsection{Baselines}
Compare Whisper zero-shot, PhoWhisper zero-shot, clean LoRA, noisy LoRA, tone-aware LoRA, and enhancement plus best ASR.

\subsection{Hardware and Training}
Local debugging uses RTX 4060 8GB and 16GB RAM. Main local training uses PhoWhisper-base, FP16, gradient checkpointing, batch size 1, gradient accumulation, max audio length 15 seconds, and LoRA rank 8. PhoWhisper-small is reserved for cloud or short ablations.

\subsection{Metrics}
Report WER, CER, syllable error rate, tone error rate, diacritic error rate, final consonant error rate, short-word deletion rate, and number/name error rate.

\section{Results and Analysis}
\subsection{Overall ASR Robustness}
Insert WER/CER by SNR table.

\subsection{Vietnamese-Specific Error Analysis}
Insert TER/DER/FCER/SWDR by SNR and noise type.

\subsection{Ablation Study}
Compare clean LoRA vs noisy LoRA, noisy LoRA vs tone-aware MTL, lambda values, and token-tone alignment policies.

\subsection{Enhancement Trade-Off}
Analyze whether enhancement lowers WER/CER but increases tone/diacritic errors.

\subsection{Failure Cases}
Show qualitative examples: tone substitution, final consonant confusion, short word deletion, names/numbers.

\section{Conclusion}
Summarize whether tone-aware multi-task learning improves Vietnamese ASR robustness under noise. State limitations: dataset scale, pseudo-label quality, token--syllable alignment, and hardware constraints. Outline future work with acoustic pitch/F0 supervision or encoder-side tone modeling.


\bibliographystyle{IEEEtran}
\bibliography{references}
\end{document}
